%% ─────────────────────────────────────────────────────────────────────────────
%% CLASE emate-ucr  —  hoja de ejercicios semanales
%% ─────────────────────────────────────────────────────────────────────────────
%%
%% Opciones de clase disponibles:
%%   [soluciones]    Muestra las soluciones dentro de cajas grises.
%%                   Sin esta opción (versión para estudiantes) las soluciones
%%                   no aparecen en el PDF, aunque estén escritas en el fuente.
%%   [numpaginas]    Siempre muestra el número de página en el pie.
%%   [nonumpaginas]  Nunca muestra el número de página.
%%   (ninguna)       Muestra el número de página solo si hay 2 o más páginas.
%%
%% Para generar ambas versiones se recomienda crear un archivo separado
%% (ver ejemplo_ejercicios_soluciones.tex):
%%
%%   \PassOptionsToClass{soluciones}{emate-ucr}
%%   %% ─────────────────────────────────────────────────────────────────────────────
%% CLASE emate-ucr  —  hoja de ejercicios semanales
%% ─────────────────────────────────────────────────────────────────────────────
%%
%% Opciones de clase disponibles:
%%   [soluciones]    Muestra las soluciones dentro de cajas grises.
%%                   Sin esta opción (versión para estudiantes) las soluciones
%%                   no aparecen en el PDF, aunque estén escritas en el fuente.
%%   [numpaginas]    Siempre muestra el número de página en el pie.
%%   [nonumpaginas]  Nunca muestra el número de página.
%%   (ninguna)       Muestra el número de página solo si hay 2 o más páginas.
%%
%% Para generar ambas versiones se recomienda crear un archivo separado
%% (ver ejemplo_ejercicios_soluciones.tex):
%%
%%   \PassOptionsToClass{soluciones}{emate-ucr}
%%   %% ─────────────────────────────────────────────────────────────────────────────
%% CLASE emate-ucr  —  hoja de ejercicios semanales
%% ─────────────────────────────────────────────────────────────────────────────
%%
%% Opciones de clase disponibles:
%%   [soluciones]    Muestra las soluciones dentro de cajas grises.
%%                   Sin esta opción (versión para estudiantes) las soluciones
%%                   no aparecen en el PDF, aunque estén escritas en el fuente.
%%   [numpaginas]    Siempre muestra el número de página en el pie.
%%   [nonumpaginas]  Nunca muestra el número de página.
%%   (ninguna)       Muestra el número de página solo si hay 2 o más páginas.
%%
%% Para generar ambas versiones se recomienda crear un archivo separado
%% (ver ejemplo_ejercicios_soluciones.tex):
%%
%%   \PassOptionsToClass{soluciones}{emate-ucr}
%%   %% ─────────────────────────────────────────────────────────────────────────────
%% CLASE emate-ucr  —  hoja de ejercicios semanales
%% ─────────────────────────────────────────────────────────────────────────────
%%
%% Opciones de clase disponibles:
%%   [soluciones]    Muestra las soluciones dentro de cajas grises.
%%                   Sin esta opción (versión para estudiantes) las soluciones
%%                   no aparecen en el PDF, aunque estén escritas en el fuente.
%%   [numpaginas]    Siempre muestra el número de página en el pie.
%%   [nonumpaginas]  Nunca muestra el número de página.
%%   (ninguna)       Muestra el número de página solo si hay 2 o más páginas.
%%
%% Para generar ambas versiones se recomienda crear un archivo separado
%% (ver ejemplo_ejercicios_soluciones.tex):
%%
%%   \PassOptionsToClass{soluciones}{emate-ucr}
%%   \input{ejemplo_ejercicios}
%%
\documentclass{emate-ucr}
% \documentclass[soluciones]{emate-ucr}   % <-- versión con soluciones

%% ── Metadatos obligatorios ───────────────────────────────────────────────────
\curso{I Ciclo 2026\\MA-1022\\Cálculo para Ciencias Económicas II} % encabezado de página
\encabezado{Ejercicios Semana N}                 % título del documento

\begin{document}

%% \imprimirtitulo
%%   Imprime el título (\encabezado) centrado en negrita.
%%   Si se definieron \fecha, \duracion o \valor (ver ejemplo_prueba_corta.tex),
%%   también muestra esa línea de metadatos bajo el título.
%%   Llamar siempre al inicio del documento, antes del primer ejercicio.
%%
\imprimirtitulo

%% ── Entorno ejercicio ────────────────────────────────────────────────────────
%%
%% \begin{ejercicio}          Ejercicio sin valor de puntos.
%% \begin{ejercicio}[10]     Ejercicio con valor indicado: muestra "(10 pts.)".
%%
%% Los ejercicios se numeran automáticamente (1, 2, 3, ...).
%% La palabra "Ejercicio" puede cambiarse con \nombreejercicio{Problema}
%% en el preámbulo (ver ejemplo_prueba_corta.tex y ejemplo_examen.tex).
%%
\begin{ejercicio}
  Resuelva cada sistema de ecuaciones lineales aplicando reducción gaussiana.
  Dé el conjunto solución.

  %% ── Entorno subejercicios ──────────────────────────────────────────────────
  %%
  %% Lista de incisos etiquetados a), b), c), ...
  %% Cada inciso va con \item.
  %%
  \begin{subejercicios}
    \item
    \[
      \begin{cases}
         x + 2y =  5 \\
        3x + 5y = 12
      \end{cases}
    \]

    %% ── Entorno solucion ────────────────────────────────────────────────────
    %%
    %% El contenido solo aparece si se cargó la opción [soluciones].
    %% Sin esa opción el bloque no genera ningún output.
    %%
    \begin{solucion}
      La matriz aumentada es
      \[
        \left(\begin{array}{cc|r} 1 & 2 & 5 \\ 3 & 5 & 12 \end{array}\right)
        \xrightarrow{f_2\to f_2-3f_1}
        \left(\begin{array}{cc|r} 1 & 2 & 5 \\ 0 & -1 & -3 \end{array}\right).
      \]
      De la fila~2: $y = 3$. De la fila~1: $x = -1$.
      El conjunto solución es $S = \{(-1,\,3)\}$.
    \end{solucion}

    \item
    \[
      \begin{cases}
         x + y = 3 \\
        2x + 2y = 8
      \end{cases}
    \]

    \begin{solucion}
      La fila~2 da $0 = 2$, contradicción. El conjunto solución es $S = \emptyset$.
    \end{solucion}

  \end{subejercicios}
\end{ejercicio}

\begin{ejercicio}[10]
  Determine el rango de la matriz
  $\begin{pmatrix} 1 & 2 & 3 \\ 2 & 5 & 7 \\ 1 & 3 & 4 \end{pmatrix}$.

  \begin{solucion}
    Reduciendo a forma escalonada se obtienen 2 filas no nulas.
    El rango es $\mathbf{2}$.
  \end{solucion}
\end{ejercicio}

\end{document}

%%
\documentclass{emate-ucr}
% \documentclass[soluciones]{emate-ucr}   % <-- versión con soluciones

%% ── Metadatos obligatorios ───────────────────────────────────────────────────
\curso{I Ciclo 2026\\MA-1022\\Cálculo para Ciencias Económicas II} % encabezado de página
\encabezado{Ejercicios Semana N}                 % título del documento

\begin{document}

%% \imprimirtitulo
%%   Imprime el título (\encabezado) centrado en negrita.
%%   Si se definieron \fecha, \duracion o \valor (ver ejemplo_prueba_corta.tex),
%%   también muestra esa línea de metadatos bajo el título.
%%   Llamar siempre al inicio del documento, antes del primer ejercicio.
%%
\imprimirtitulo

%% ── Entorno ejercicio ────────────────────────────────────────────────────────
%%
%% \begin{ejercicio}          Ejercicio sin valor de puntos.
%% \begin{ejercicio}[10]     Ejercicio con valor indicado: muestra "(10 pts.)".
%%
%% Los ejercicios se numeran automáticamente (1, 2, 3, ...).
%% La palabra "Ejercicio" puede cambiarse con \nombreejercicio{Problema}
%% en el preámbulo (ver ejemplo_prueba_corta.tex y ejemplo_examen.tex).
%%
\begin{ejercicio}
  Resuelva cada sistema de ecuaciones lineales aplicando reducción gaussiana.
  Dé el conjunto solución.

  %% ── Entorno subejercicios ──────────────────────────────────────────────────
  %%
  %% Lista de incisos etiquetados a), b), c), ...
  %% Cada inciso va con \item.
  %%
  \begin{subejercicios}
    \item
    \[
      \begin{cases}
         x + 2y =  5 \\
        3x + 5y = 12
      \end{cases}
    \]

    %% ── Entorno solucion ────────────────────────────────────────────────────
    %%
    %% El contenido solo aparece si se cargó la opción [soluciones].
    %% Sin esa opción el bloque no genera ningún output.
    %%
    \begin{solucion}
      La matriz aumentada es
      \[
        \left(\begin{array}{cc|r} 1 & 2 & 5 \\ 3 & 5 & 12 \end{array}\right)
        \xrightarrow{f_2\to f_2-3f_1}
        \left(\begin{array}{cc|r} 1 & 2 & 5 \\ 0 & -1 & -3 \end{array}\right).
      \]
      De la fila~2: $y = 3$. De la fila~1: $x = -1$.
      El conjunto solución es $S = \{(-1,\,3)\}$.
    \end{solucion}

    \item
    \[
      \begin{cases}
         x + y = 3 \\
        2x + 2y = 8
      \end{cases}
    \]

    \begin{solucion}
      La fila~2 da $0 = 2$, contradicción. El conjunto solución es $S = \emptyset$.
    \end{solucion}

  \end{subejercicios}
\end{ejercicio}

\begin{ejercicio}[10]
  Determine el rango de la matriz
  $\begin{pmatrix} 1 & 2 & 3 \\ 2 & 5 & 7 \\ 1 & 3 & 4 \end{pmatrix}$.

  \begin{solucion}
    Reduciendo a forma escalonada se obtienen 2 filas no nulas.
    El rango es $\mathbf{2}$.
  \end{solucion}
\end{ejercicio}

\end{document}

%%
\documentclass{emate-ucr}
% \documentclass[soluciones]{emate-ucr}   % <-- versión con soluciones

%% ── Metadatos obligatorios ───────────────────────────────────────────────────
\curso{I Ciclo 2026\\MA-1022\\Cálculo para Ciencias Económicas II} % encabezado de página
\encabezado{Ejercicios Semana N}                 % título del documento

\begin{document}

%% \imprimirtitulo
%%   Imprime el título (\encabezado) centrado en negrita.
%%   Si se definieron \fecha, \duracion o \valor (ver ejemplo_prueba_corta.tex),
%%   también muestra esa línea de metadatos bajo el título.
%%   Llamar siempre al inicio del documento, antes del primer ejercicio.
%%
\imprimirtitulo

%% ── Entorno ejercicio ────────────────────────────────────────────────────────
%%
%% \begin{ejercicio}          Ejercicio sin valor de puntos.
%% \begin{ejercicio}[10]     Ejercicio con valor indicado: muestra "(10 pts.)".
%%
%% Los ejercicios se numeran automáticamente (1, 2, 3, ...).
%% La palabra "Ejercicio" puede cambiarse con \nombreejercicio{Problema}
%% en el preámbulo (ver ejemplo_prueba_corta.tex y ejemplo_examen.tex).
%%
\begin{ejercicio}
  Resuelva cada sistema de ecuaciones lineales aplicando reducción gaussiana.
  Dé el conjunto solución.

  %% ── Entorno subejercicios ──────────────────────────────────────────────────
  %%
  %% Lista de incisos etiquetados a), b), c), ...
  %% Cada inciso va con \item.
  %%
  \begin{subejercicios}
    \item
    \[
      \begin{cases}
         x + 2y =  5 \\
        3x + 5y = 12
      \end{cases}
    \]

    %% ── Entorno solucion ────────────────────────────────────────────────────
    %%
    %% El contenido solo aparece si se cargó la opción [soluciones].
    %% Sin esa opción el bloque no genera ningún output.
    %%
    \begin{solucion}
      La matriz aumentada es
      \[
        \left(\begin{array}{cc|r} 1 & 2 & 5 \\ 3 & 5 & 12 \end{array}\right)
        \xrightarrow{f_2\to f_2-3f_1}
        \left(\begin{array}{cc|r} 1 & 2 & 5 \\ 0 & -1 & -3 \end{array}\right).
      \]
      De la fila~2: $y = 3$. De la fila~1: $x = -1$.
      El conjunto solución es $S = \{(-1,\,3)\}$.
    \end{solucion}

    \item
    \[
      \begin{cases}
         x + y = 3 \\
        2x + 2y = 8
      \end{cases}
    \]

    \begin{solucion}
      La fila~2 da $0 = 2$, contradicción. El conjunto solución es $S = \emptyset$.
    \end{solucion}

  \end{subejercicios}
\end{ejercicio}

\begin{ejercicio}[10]
  Determine el rango de la matriz
  $\begin{pmatrix} 1 & 2 & 3 \\ 2 & 5 & 7 \\ 1 & 3 & 4 \end{pmatrix}$.

  \begin{solucion}
    Reduciendo a forma escalonada se obtienen 2 filas no nulas.
    El rango es $\mathbf{2}$.
  \end{solucion}
\end{ejercicio}

\end{document}

%%
\documentclass{emate-ucr}
% \documentclass[soluciones]{emate-ucr}   % <-- versión con soluciones

%% ── Metadatos obligatorios ───────────────────────────────────────────────────
\curso{I Ciclo 2026\\MA-1022\\Cálculo para Ciencias Económicas II} % encabezado de página
\encabezado{Ejercicios Semana N}                 % título del documento

\begin{document}

%% \imprimirtitulo
%%   Imprime el título (\encabezado) centrado en negrita.
%%   Si se definieron \fecha, \duracion o \valor (ver ejemplo_prueba_corta.tex),
%%   también muestra esa línea de metadatos bajo el título.
%%   Llamar siempre al inicio del documento, antes del primer ejercicio.
%%
\imprimirtitulo

%% ── Entorno ejercicio ────────────────────────────────────────────────────────
%%
%% \begin{ejercicio}          Ejercicio sin valor de puntos.
%% \begin{ejercicio}[10]     Ejercicio con valor indicado: muestra "(10 pts.)".
%%
%% Los ejercicios se numeran automáticamente (1, 2, 3, ...).
%% La palabra "Ejercicio" puede cambiarse con \nombreejercicio{Problema}
%% en el preámbulo (ver ejemplo_prueba_corta.tex y ejemplo_examen.tex).
%%
\begin{ejercicio}
  Resuelva cada sistema de ecuaciones lineales aplicando reducción gaussiana.
  Dé el conjunto solución.

  %% ── Entorno subejercicios ──────────────────────────────────────────────────
  %%
  %% Lista de incisos etiquetados a), b), c), ...
  %% Cada inciso va con \item.
  %%
  \begin{subejercicios}
    \item
    \[
      \begin{cases}
         x + 2y =  5 \\
        3x + 5y = 12
      \end{cases}
    \]

    %% ── Entorno solucion ────────────────────────────────────────────────────
    %%
    %% El contenido solo aparece si se cargó la opción [soluciones].
    %% Sin esa opción el bloque no genera ningún output.
    %%
    \begin{solucion}
      La matriz aumentada es
      \[
        \left(\begin{array}{cc|r} 1 & 2 & 5 \\ 3 & 5 & 12 \end{array}\right)
        \xrightarrow{f_2\to f_2-3f_1}
        \left(\begin{array}{cc|r} 1 & 2 & 5 \\ 0 & -1 & -3 \end{array}\right).
      \]
      De la fila~2: $y = 3$. De la fila~1: $x = -1$.
      El conjunto solución es $S = \{(-1,\,3)\}$.
    \end{solucion}

    \item
    \[
      \begin{cases}
         x + y = 3 \\
        2x + 2y = 8
      \end{cases}
    \]

    \begin{solucion}
      La fila~2 da $0 = 2$, contradicción. El conjunto solución es $S = \emptyset$.
    \end{solucion}

  \end{subejercicios}
\end{ejercicio}

\begin{ejercicio}[10]
  Determine el rango de la matriz
  $\begin{pmatrix} 1 & 2 & 3 \\ 2 & 5 & 7 \\ 1 & 3 & 4 \end{pmatrix}$.

  \begin{solucion}
    Reduciendo a forma escalonada se obtienen 2 filas no nulas.
    El rango es $\mathbf{2}$.
  \end{solucion}
\end{ejercicio}

\end{document}
